\subsection{A missing wild semistability defect in characteristic three}

\begin{theorem}[B-RD8: defect twelve with full rational $3$-torsion]\label{thm:B-RD8}
\status{Proved; exact computation recorded.}
The list $2,3,4,6$ in B.2 is incomplete in residue characteristic three.
Put
\[
 K=\mathbf Q_3(i,\zeta),\qquad i^2=-1,\quad
 \zeta^2+\zeta+1=0,\qquad \pi=\zeta-1,\quad t=1+i\pi.
\]
Then $K$ has absolute ramification index $2$ and residue field
$\mathbf F_9$. The elliptic curve
\[
 E:\quad y^2=x^3-27t(t^3+8)x+54(t^6-20t^3-8)
\]
has full $3$-torsion over $K$ and additive potentially supersingular
reduction, with semistability defect exactly $12$.
\end{theorem}
\begin{proof}
Write $a=-27t(t^3+8)$ and $b=54(t^6-20t^3-8)$.
Set $\delta=1+2\zeta=3+2\pi$, so $\delta^2=-3$.
The eight nonzero $3$-torsion points are
\[
 \left(-9t^2,\ \pm12\delta(t^3-1)\right)
\]
and, for $\epsilon\in\{1,\zeta,\zeta^2\}$,
\[
 \left(3t^2+12\epsilon t+12\epsilon^2,
       \ \pm36(\epsilon t^2+\epsilon^2t+1)\right).
\]
Indeed their four $x$-coordinates give the factorization
\[
 3x^4+6ax^2+12bx-a^2
 =3(x+9t^2)\prod_{\epsilon^3=1}
       (x-3t^2-12\epsilon t-12\epsilon^2).
\]
Substitution verifies their $y$-coordinates; the points are distinct
and exhaust the eight nonzero geometric $3$-torsion points.
The exact certificate \texttt{verify\_p3\_defect12.py} performs these
calculations in
$\mathbf Q[i,\pi]/(i^2+1,\pi^2+3\pi+3)$, verifies the curve equation
and the tripling identity for every point, and verifies independence
of points with different $x$-coordinates. No numerical approximation
is used.

Normalize $v_K(\pi)=1$. Direct calculation gives
\[
 \Delta=2^{12}3^9(t^3-1)^3,\quad
 v_K(t^3-1)=v_K(t^3+8)=3,\quad v_K(\Delta)=27.
\]
In particular $E$ is not good over $K$, since the valuation of a
Weierstrass discriminant can change only by a multiple of $12$.
We next both construct good reduction and prove the required minimality.

In the $2$-division polynomial $x^3+ax+b$, substitute
$x=3\pi^2 Z$. Dividing by $(3\pi^2)^3$ gives
\[
 Z^3+A Z+B,
 \qquad A=4i\pi+3\pi+6,\quad B=-12i+20\pi+30.
\]
Both $A$ and $B$ have valuation one. The polynomial is Eisenstein,
so a root $z$ generates a totally ramified cubic extension $K_1/K$,
and $v_K(z)=1/3$. Set $x_0=3\pi^2 z$. Then
\[
 v_K(x_0)=13/3,\qquad v_K(3x_0^2+a)=9.
\]
The second equality can also be read from the derivative
$3z^2+A$, whose valuation is one. Adjoin $u$ satisfying
\[
 u^4=3x_0^2+a.
\]
In the normalized valuation of $K_1$, the right side has value $27$,
prime to four. Thus this is a totally ramified degree-four extension;
for $L=K_1(u)$ we have $e(L/K)=12$ and $v_K(u)=9/4$.
The substitution $x=x_0+u^2X$, $y=u^3Y$ gives
\[
 Y^2=X^3+\frac{3x_0}{u^2}X^2+X.
\]
The middle coefficient has valuation $11/6>0$. The special fibre is
$Y^2=X^3+X$, smooth and supersingular in characteristic three.
Therefore good reduction is obtained with ramification degree twelve.

Conversely, let $F/K$ be any finite extension over which $E$ has good
reduction. Its discriminant valuation implies
$12\mid27e(F/K)$, hence $4\mid e(F/K)$. Since $2$ is invertible,
the $2$-torsion of a good elliptic scheme is finite etale; all its
points are defined over the maximal unramified extension of $F$.
The Eisenstein $2$-division root has $K$-valuation $1/3$ after the
indicated scaling, so $3\mid e(F/K)$. Consequently
$12\mid e(F/K)$. Together with the constructed extension, this proves
that the semistability defect is exactly twelve.
\end{proof}
